\section{講演を聴講しての感想}
Transformerが転機となったことは理解していたが，講演を通して，その後に国内でLLM.jpのような組織が立ち上がる流れを把握できた．大規模言語モデルは地域固有の知識を十分に取り込む必要があり，汎用モデルだけでは不十分であること，さらに地政学的リスクへの備えとして先端技術を自国でも扱える体制づくりが重要であるという指摘が印象的だった．最先端そのものを追う必要はなくとも，「自分たちで動かせる最新技術」を確保しておくことの重要性が語られた．

日本語は2文字で1トークンになるため，数兆規模の学習データが必要となる一方，ネット上で得られる日本語コーパスは約1兆トークンに留まるという現実も示された．この状況が，国産LLM開発の難しさと必要性の双方を裏付けている．

LLM.jpは，国産モデルをつくるだけでなく，オープンソース開発やデータ公開を中心的な活動としており，ライセンスはApatcheを採用している．研究現場ではLLMを使う場面が増えているが，日本語に強いモデルが不可欠なケースは多く，こうした基盤を国内で整備する意義は大きいと感じた．